\subsection{Zweck der Deliver-Phase}

Die Deliver-Phase bildet den zweiten Teil des zweiten Diamanten und überführt die in der Develop-Phase entwickelten Prototypen in eine validierte Stossrichtung \parencite{designcouncil2005}. Während die Develop-Phase den Lösungsraum geöffnet und das Konzept \enquote{ZVV Tageserlebnis} erfahrbar gemacht hat, beantwortet die Deliver-Phase die Frage \enquote{Trägt die Lösung wirklich?}. Im Zentrum steht damit das Testen der Annahmen, auf denen das Konzept beruht: Akzeptanz, Preisbereitschaft und digitale Hürden.

Methodisch folgt die Phase dem Grundsatz, einen Prototyp zu bauen, als hätte man recht, ihn aber zu testen, als läge man falsch \parencite{leifer2015}. Nicht die Bestätigung der eigenen Idee ist das Ziel, sondern das frühe Aufdecken von Schwachstellen, solange Anpassungen noch günstig sind. Die Prototypen aus der Develop-Phase konfrontieren wir dafür mit echten Nutzenden aus der Zielgruppe und werten sie systematisch aus \parencite{zhaw2026testing}.

\subsection{Eingesetzte Methoden und Vorgehen}

\subsubsection{Qualitativer Prototypentest}
Zur Validierung führen wir am 19.~Juni 2026 zwei qualitative Prototypentests durch. Getestet werden drei Prototypen des Gesamtkonzepts: die \textbf{Planer-App} für den Gastgeber, die \textbf{Anmelde-App} für die eingeladenen Teilnehmenden sowie die \textbf{Treffpunkt-Idee \enquote{Donnerstags mit Kurt}}. Jeder Test folgt einem strukturierten \gls{leitfadeninterview} mit einem festen Ablauf: einer Einführung ins Projekt, einem kurzen Video zur Einstimmung, anschliessend je einer kurzen Demo der beiden Apps und zuletzt der Vorstellung der Treffpunkt-Idee über Plakat und Website. Zu jedem Prototyp stellen wir dieselbe Abfolge offener Fragen: ob die Person die Lösung nutzen würde, was ihr gefällt, was nicht und welche Alternativen ihr einfallen.

\subsubsection{Auswahl der Testpersonen}
Als Testpersonen wählen wir bewusst zwei pensionierte ehemalige Lehrpersonen aus, \textbf{Lisbeth} (68) und \textbf{Enrico}, genannt Rico (71). Der Auswahl liegt die Annahme zugrunde, dass sich Lehrpersonen im Berufsleben überdurchschnittlich häufig mit der Organisation von Gruppenreisen befassen, etwa Klassenlagern, Exkursionen und Schulreisen, und daher der Gastgeber- und Planerrolle der Persona Kurt besonders nahe stehen. Zugleich decken die beiden zwei gegensätzliche Reisetypen ab: Lisbeth ist reisefreudig, ein ausgesprochener Gruppenmensch und plant gerne für ihre Familie, während Rico ein Einzelgänger ist, am liebsten allein reist und den \gls{oev} mangels Halbtax als zu teuer empfindet. So lassen sich sowohl die Sicht potenzieller Gastgeber als auch die kritische Aussenperspektive abbilden.

\subsubsection{Feedback Capture Grid}
Zur Auswertung nutzen wir je Prototyp ein \gls{feedbackgrid}. Dieses ordnet alle Rückmeldungen in vier Quadranten, \textbf{Positiv}, \textbf{Optimieren}, \textbf{Fragen} und \textbf{Ideen}, und macht die Beobachtungen so über beide Testpersonen hinweg vergleichbar \parencite{zhaw2026testing}. Aussagekräftige Originalzitate ordnen wir den jeweiligen Personen zu, um die Befunde nachvollziehbar zu belegen.

\newpage
\subsection{Zentrale inhaltliche Ergebnisse}

\subsubsection{Prototyp 1: Planer-App}
Die Planer-App stösst bei beiden Testpersonen auf sofortige Begeisterung: Den Aufbau erleben sie als klar und verständlich, die App-Idee erfassen sie schnell. Besonders treffend ist Lisbeths Vergleich, die App sei \enquote{wie ein Doodle, aber für eine Reise}, bei der die Teilnehmenden alles auf einen Blick sehen. Ausdrücklich loben beide die Erinnerungsfunktion und die Übersicht, wer sich angemeldet hat, den sichtbaren Preis als wichtiges Informationselement sowie die Wissens- und Geschichtsinhalte der Touren als erkennbaren Mehrwert (\enquote{Überall ist Wissen, das finde ich gut}, Rico). Bemerkenswert ist Ricos Einordnung der Kostenfrage: \enquote{Wenn so ein Erlebnis dahinter ist, würde ich bezahlen}, der \gls{oev}-Preis ist kein Hindernis, sofern das Erlebnis stimmt.

Im Quadranten \emph{Optimieren} zeigen sich konkrete Verbesserungswünsche: eine flexibel wählbare Startzeit (nicht jeder kann um 09:00~Uhr starten), ein flexibler Einstiegsbahnhof (nicht alle kommen aus Zürich), eine begrenzbare Teilnehmerzahl von rund 15 Personen, ein klarer Hinweis auf den Preis für Vollzahler ohne Halbtax sowie eine deutlichere Beschreibung des Ticket-Inhalts. Offen bleiben Fragen nach der Angabe von Verpflegungswünschen, nach dem Zugang ohne Smartphone und danach, ob direkt in der App gebucht oder nur geplant wird.

\subsubsection{Prototyp 2: Anmelde-App}
Die Anmelde-App aus Sicht der eingeladenen Teilnehmenden überzeugt durch ihre Übersichtlichkeit: \enquote{Sehr gut, gute Übersicht, man weiss genau, was man da macht} (Lisbeth). Lisbeth würde die App nutzen und sie als Vereinfachung des Einladens von Freunden und Familie schätzen. Auch Rico, der selbst lieber allein reist, erkennt den Wert für sein Umfeld klar an: \enquote{In meinem Umfeld würden es viele machen, das weiss ich}.

Im Zentrum der Optimierungspunkte steht der Bezahlprozess, der für beide unklar bleibt (\enquote{Wie zahle ich per Twint? Mir ist nicht klar, wie bezahlt wird}, Lisbeth), ergänzt um den erneuten Wunsch nach einer klar ausgewiesenen Halbtax-Variante und einem verständlicheren Ticket-Inhalt. Aus den offenen Fragen ragt eine besonders heraus: ob Kurt am Treffpunkt mit dem Geldeinzug beschäftigt sei, \enquote{Würde ich mir fragen, ob Kurt die erste halbe Stunde damit beschäftigt ist, Geld einzuziehen} (Lisbeth). Daraus entsteht die übereinstimmende Idee beider, die Bezahlung vollständig vorab und digital zu lösen, kein Bargeld am Treffpunkt, sowie Verpflegungswünsche bereits bei der Anmeldung zu erfassen.

\subsubsection{Prototyp 3: Treffpunkt \enquote{Donnerstags mit Kurt}}
Der wöchentliche Offline-Treffpunkt am Zürcher Hauptbahnhof ist der klare Favorit beider Testpersonen. Das Konzept verstehen beide sofort und bewerten es als gute Idee, Rico sieht ein deutliches Potenzial gerade für Pensionierte, die allein oder frisch pensioniert sind. Der niederschwellige, gesellige Charakter überzeugt, weil er Menschen ausserhalb des eigenen Umfelds zusammenbringt.

Zugleich nennen beide klare Bedingungen: Das Programm müsse vorab transparent sein (\enquote{nicht hingehen, ohne zu wissen wohin}), die Bezahlung müsse geklärt sein, damit Kurt nicht am Treffpunkt Geld einziehen muss, und statt eines QR-Scans brauche es einen einfachen Klick-Link. Damit bestätigt der dritte Prototyp die übergreifende Tendenz: Der soziale, physische Anker trägt am stärksten, sofern Information und Bezahlung im Vorfeld gelöst sind.

\subsubsection{Vier zentrale Erkenntnisse}
Über alle drei Prototypen hinweg verdichten sich die Beobachtungen zu vier Erkenntnissen:

\begin{description}
  \item[1. Vorinformation ist Pflicht.] Niemand fährt blind los. Ziel, Programm, Treffpunkt, Preis und Verpflegung müssen vorab klar kommuniziert sein.
  \item[2. Der Treffpunkt zündet am stärksten.] Das niederschwellige Donnerstags-Format spricht Pensionierte am meisten an und bringt Fremde zusammen.
  \item[3. Digital darf keine Hürde sein.] QR-Code, App und Online-Zahlung überfordern viele in der Zielgruppe. Es braucht den einfachen Link und stets auch einen analogen Weg.
  \item[4. Kurt ist der Schlüssel, entlastet ihn.] Der vertraute Gastgeber macht den Unterschied. Tickets und Bezahlung dürfen ihn aber nicht aufhalten, und seine Spesen müssen gedeckt sein.
\end{description}

\subsubsection{Mehrwert für den ZVV}
Aus den bestätigten Annahmen lässt sich der Nutzen für den \gls{zvv} klar benennen. Tagesausflüge füllen \textbf{freie Sitze in den Nebenzeiten} unter der Woche und ausserhalb des Pendelverkehrs, ohne zusätzliche Infrastruktur-Investition. Das Angebot erschliesst eine \textbf{neue Zielgruppe 65+}, die sonst selten den \gls{oev} nutzt, und bindet sie als regelmässige Kundschaft. Es stiftet zugleich einen \textbf{gesellschaftlichen Wert}, indem es gegen Vereinsamung wirkt und ältere Menschen aktiv hält, was direkt auf die Marke einzahlt. Und es hat eine \textbf{niedrige Einstiegshürde}, weil es auf dem bestehenden Netz und Tarif aufbaut und damit in rund drei Monaten pilotfähig wäre.

\subsubsection{Nächste Schritte}
Aus den Tests leiten wir einen vierstufigen Weg vom Prototyp zum Pilotbetrieb ab:

\begin{enumerate}
  \item \textbf{Pilot mit drei Touren} in einer Region mit ausgewählten Gastgebern und echtem Feedback der Zielgruppe.
  \item \textbf{Integration des \gls{oev}-Tickets}, also die saubere Anbindung des Festpreis-Pakets an Tarif und Buchung des \gls{zvv}.
  \item \textbf{Touren ausbauen} mit weiteren Themen, Bausteinen und Regionen.
  \item \textbf{Rollout im Kanton} als festes Angebot mit vielen Gastgebern und vollen Wochentagen.
\end{enumerate}

\subsection{Dokumentation der Feedback Capture Grids}
Die folgenden drei Raster dokumentieren die Rohrückmeldungen beider Testpersonen je Prototyp, geordnet nach den vier Quadranten des \gls{feedbackgrid}. Sie bilden die Grundlage der oben verdichteten Ergebnisse. Kürzel: \textbf{L}~=~Lisbeth, \textbf{R}~=~Rico, \textbf{Beide}~=~übereinstimmende Aussage.

\end{multicols}

\newpage
\begin{center}\textbf{Prototyp 1: Tagestouren-App (Planer-App mit fertig zusammengestellten Ausflügen)}\end{center}
\FeedbackGrid
{% Positiv
\begin{itemize}[nosep,leftmargin=1.3em]
  \item Sofortige Begeisterung, klare Übersicht, Idee schnell erfasst \textbf{(L\,\&\,R)}
  \item \enquote{Wie ein Doodle, aber für eine Reise}, alles auf einen Blick \textbf{(L)}
  \item Sichtbarer Preis als wichtiges Info-Element \textbf{(L)}
  \item Wissensinhalte als Mehrwert: \enquote{Überall ist Wissen} \textbf{(R)}
  \item ÖV-Preis kein Hindernis: \enquote{Mit Erlebnis dabei, würde ich bezahlen} \textbf{(R)}
\end{itemize}}
{% Optimieren
\begin{itemize}[nosep,leftmargin=1.3em]
  \item Startzeit flexibel wählbar, nicht jeder startet um 09:00 \textbf{(L)}
  \item Einstiegsbahnhof flexibel, nicht alle kommen aus Zürich \textbf{(R)}
  \item Datumsvorschläge im Doodle-Stil anbieten \textbf{(L)}
  \item Teilnehmerzahl limitierbar, max.~15 Personen \textbf{(R)}
  \item Ticket-Inhalt klarer: was ist inklusive (Essen, Eintritt)? \textbf{(R)}
\end{itemize}}
{% Fragen
\begin{itemize}[nosep,leftmargin=1.3em]
  \item Verpflegungswünsche angebbar (vegetarisch, kein Mittagessen)? \textbf{(R)}
  \item Zugang zum Ticket ohne Smartphone? \textbf{(L)}
  \item Tour passend zur Fitness der Gruppe wählen? \textbf{(R)}
  \item Direkt in der App gebucht oder nur geplant? \textbf{(L)}
\end{itemize}}
{% Ideen
\begin{itemize}[nosep,leftmargin=1.3em]
  \item Mitteilungsfeld für individuelle Wünsche bei der Anmeldung \textbf{(L)}
  \item Routen anpassbar: Einstieg ab beliebiger Haltestelle \textbf{(L)}
  \item Mehrere Abfahrtszeiten pro Tour als Optionen \textbf{(L)}
  \item Touren auch für Haustiere (Hunde), Nischenidee \textbf{(L)}
\end{itemize}}
\captionof{figure}{Feedback Capture Grid, Prototyp~1: Tagestouren-App}

\bigskip
\begin{center}\textbf{Prototyp 2: Gruppenreise-App (Teilnehmer-Ansicht und Buchungsflow)}\end{center}
\FeedbackGrid
{% Positiv
\begin{itemize}[nosep,leftmargin=1.3em]
  \item Gute Übersicht: \enquote{man weiss genau, was man da macht} \textbf{(L)}
  \item Vereinfacht das Einladen von Freunden und Familie, \enquote{würde es auf jeden Fall nutzen} \textbf{(L)}
  \item Wert fürs Umfeld anerkannt: \enquote{In meinem Umfeld würden es viele machen} \textbf{(R)}
\end{itemize}}
{% Optimieren
\begin{itemize}[nosep,leftmargin=1.3em]
  \item Bezahlprozess unklar: \enquote{Wie zahle ich per Twint?} \textbf{(L)}
  \item Ticket-Inhalt klarer kommunizieren, ein Sternchen reicht nicht \textbf{(R)}
  \item Halbtax-Variante separat ausweisen \textbf{(L)}
\end{itemize}}
{% Fragen
\begin{itemize}[nosep,leftmargin=1.3em]
  \item Ist Kurt am HB mit Geldeinzug beschäftigt? \enquote{\dots die erste halbe Stunde Geld einziehen} \textbf{(L)}
  \item Wie wird die Buchung bestätigt, Ticket per SMS oder selbst lösen? \textbf{(L)}
\end{itemize}}
{% Ideen
\begin{itemize}[nosep,leftmargin=1.3em]
  \item Verpflegungswünsche bereits bei der Anmeldung erfassen \textbf{(L)}
  \item Bezahlung komplett vorab digital, kein Bargeld am Treffpunkt \textbf{(Beide)}
  \item Maximale Gruppengrösse 12--15, danach ein zweiter Guide \textbf{(R)}
\end{itemize}}
\captionof{figure}{Feedback Capture Grid, Prototyp~2: Gruppenreise-App}

\bigskip
\begin{center}\textbf{Prototyp 3: \enquote{Donnerstags mit Kurt}, Offline-Treffpunkt Zürich HB}\end{center}
\FeedbackGrid
{% Positiv
\begin{itemize}[nosep,leftmargin=1.3em]
  \item Konzept klar verstanden, gute Idee, \enquote{gerade für aktive Rentner} \textbf{(L)}
  \item Potenzial für allein oder frisch Pensionierte: \enquote{wie Wandergruppen, einfach mit ÖV} \textbf{(R)}
  \item Offenheit für Fremde: \enquote{mitgehen, um neue Leute kennenzulernen} \textbf{(L)}
  \item HB und Kiosk als natürlicher Sammelpunkt bestätigt \textbf{(L\,\&\,R)}
  \item Rico würde selbst ein \enquote{Kurt} werden, max.~12, Reise gratis \textbf{(R)}
  \item Erwartung an Kurt: Wissen, Geschichten erzählen, Lernen ermöglichen \textbf{(L)}
\end{itemize}}
{% Optimieren
\begin{itemize}[nosep,leftmargin=1.3em]
  \item Vorabinformation zwingend, nicht \enquote{blind} auftauchen \textbf{(L)}
  \item Bezahlprozess vorab klären \textbf{(L)}
  \item QR-Code wird nicht gescannt~$\rightarrow$~einfacher Weblink/URL \textbf{(R)}
  \item Werbung gezielt platzieren: Facebook, WhatsApp, Gemeindeanlässe \textbf{(R)}
  \item Recruiting-Film auf Mundart sprechen \textbf{(R)}
  \item Video länger und erklärender gestalten \textbf{(R)}
\end{itemize}}
{% Fragen
\begin{itemize}[nosep,leftmargin=1.3em]
  \item Bezahlung am Treffpunkt: Twint, Cash oder vorab? \textbf{(L)}
  \item Wie wird man offiziell ein \enquote{Kurt}, Anmeldung, Vorbereitung? \textbf{(R)}
  \item Wie viele Donnerstags-Ziele pro Monat, wie weit voraus sichtbar? \textbf{(L)}
  \item Treffpunkt auch in anderen Städten (Luzern, Bern)? \textbf{(L)}
\end{itemize}}
{% Ideen
\begin{itemize}[nosep,leftmargin=1.3em]
  \item Online-Kalender mit allen Donnerstags-Zielen (ohne App) \textbf{(L)}
  \item Facebook und WhatsApp als Hauptkanal für 65+ \textbf{(R)}
  \item Distribution über Gemeinde-Newsletter und Rentnerveranstaltungen \textbf{(R)}
  \item Ausweitung auf Luzern, Bern, Winterthur \textbf{(L)}
\end{itemize}}
\captionof{figure}{Feedback Capture Grid, Prototyp~3: \enquote{Donnerstags mit Kurt}}

\begin{multicols}{2}

\subsection{Reflexion der Methode}

Der Prototypentest bestätigt den Wert des testorientierten Mindsets: Erst die Konfrontation mit echten Nutzenden macht sichtbar, dass nicht die Apps, sondern der physische Treffpunkt am stärksten trägt, und dass die digitale Bezahlung eine ernsthafte Hürde darstellt. Beide Befunde sind aus der reinen Konzeptarbeit nicht ableitbar und korrigieren die ursprüngliche, app-zentrierte Erwartung. Das Feedback Capture Grid erweist sich dabei als wirksames Werkzeug, um die Fülle an Rückmeldungen strukturiert und vergleichbar zu verdichten.

Kritisch einzuordnen ist die geringe Stichprobe von nur zwei qualitativen Tests, die keine repräsentativen Aussagen erlaubt und erste Tendenzen statt belastbarer Häufigkeiten liefert. Hinzu kommt, dass die Prototypen teils nur telefonisch gezeigt werden können, was die Demo einschränkt, und dass die Auswahl ehemaliger Lehrpersonen auf einer plausiblen, aber nicht validierten Annahme über deren Nähe zur Gastgeberrolle beruht. Die Tests ersetzen daher keinen Feldtest, sondern schärfen die Richtung: Die abgeleiteten nächsten Schritte überführen die offenen Punkte, allen voran die Vorab-Bezahlung und der niederschwellige Zugang, in einen echten Pilotbetrieb, in dem sie mit einer grösseren und breiteren Gruppe validiert werden.
